\documentclass[11pt]{article}
\usepackage[T1]{fontenc}
\usepackage[utf8]{inputenc}
\usepackage{lmodern,amsmath,amssymb,amsthm,mathtools,booktabs,microtype}
\usepackage[margin=1in]{geometry}
\usepackage[colorlinks=true,linkcolor=blue,citecolor=blue,urlcolor=blue]{hyperref}
\hypersetup{pdftitle={Irrational even beta values from powered hypergeometric kernels},
pdfauthor={twinheterodox},
pdfsubject={Working proof draft with exact computational certificates}}
\usepackage{enumitem}
\setlist{nosep}
\newtheorem{theorem}{Theorem}
\newtheorem{proposition}[theorem]{Proposition}
\newtheorem{lemma}[theorem]{Lemma}
\newtheorem{corollary}[theorem]{Corollary}
\newcommand{\Z}{\mathbb Z}\newcommand{\Q}{\mathbb Q}
\newcommand{\Fp}{\mathbb F_p}\newcommand{\vp}{v_p}
\newcommand{\ind}{\mathbf 1}
\title{Irrational even beta values from powered hypergeometric kernels}
\author{twinheterodox\\[2pt]\normalsize Prepared with Codex}
\date{11 November 2026}
\begin{document}
\maketitle

\begin{abstract}
For the Dirichlet beta function, we prove that at least one of
$\beta(6),\beta(8),\ldots,\beta(74)$ is irrational, and that at least one
of $\beta(8),\beta(10),\ldots,\beta(110)$ is irrational. These shorten the
previously stated intervals with the same starting indices, whose upper
endpoints are $92$ and $128$, respectively \cite[Section 8.2]{shifted}.
The new bounds come from a different construction, based on powers of
Krisman's hypergeometric kernel \cite{krisman}.
\end{abstract}

\section{Results and proof strategy}\label{sec:introduction}
For $s>1$, the Dirichlet beta function is
\[
 \beta(s)=\sum_{t=0}^{\infty}\frac{(-1)^t}{(2t+1)^s}.
\]
Krisman \cite[Theorem 1.1]{krisman} proves that at least one of
$\beta(4),\beta(6),\ldots,\beta(36)$ is irrational. We seek intervals
that also omit $\beta(4)$, and in the second case below omit $\beta(6)$.

\begin{theorem}\label{thm:main}
Each of the following sets contains an irrational number:
\[
 \{\beta(6),\beta(8),\ldots,\beta(74)\},\qquad
 \{\beta(8),\beta(10),\ldots,\beta(110)\}.
\]
\end{theorem}

Table~\ref{tab:comparison} compares these bounds with the earlier
shifted-interval results \cite{shifted}. Each entry asserts the existence of
one irrational value, with all indices even. The intervals overlap, so
the two assertions of Theorem~\ref{thm:main} need not identify distinct
irrational numbers.
\begin{table}[ht]
\centering
\begin{tabular}{crrrr}\toprule
 & \multicolumn{2}{c}{\cite[Section 8.2]{shifted}} &
   \multicolumn{2}{c}{Theorem~\ref{thm:main}}\\
First index & Last index & Values & Last index & Values\\\midrule
6 & 92 & 44 & 74 & 35\\
8 & 128 & 61 & 110 & 52\\\bottomrule
\end{tabular}
\caption{The earlier endpoints use parameters chosen for these particular
starting indices. Each new interval removes nine candidate values.}
\label{tab:comparison}
\end{table}

Sorokin \cite[Theorem S]{sorokin} proved that, for every sufficiently
large odd $m$ and every even $h>100m\log m$, at least one of
$\beta(m+1),\beta(m+3),\ldots,\beta(m+h-1)$ is irrational.
The later result \cite[Theorem 1.2]{shifted} treats every even
starting index and any prescribed number of irrational values. For one
irrational value, \cite[Proposition B.1]{shifted} gives an upper endpoint below $18.58m$
for positive odd $m$, with first index $m+1$. We improve the two instances
in the table. Corollary~\ref{cor:general} gives a criterion for further
instances, with a finite nonvanishing condition that must still be checked.

\subsection*{How the argument is organized}
The irrationality contradiction has three requirements. For each target
interval, we construct a sequence of rational linear forms in $1$ and
the indicated beta values. After clearing the denominators of their
coefficients, these forms must tend to zero and must be nonzero infinitely
often. Under joint
rationality, a single fixed denominator for the beta values would then
give a sequence of integers tending to zero, with infinitely many nonzero
terms.

Krisman's construction already balances the cost of clearing denominators
against the decay of a rational kernel. We retain that balance while
changing the beta indices. Squaring the kernel gives zeros of order six,
so four differentiations still leave zeros at the integers where they are
needed. Reflection then makes $\beta(6)$ the first possible beta term.
Cubing and differentiating six times similarly starts at $\beta(8)$.
Section~\ref{sec:construction} explains this operation before writing its
general form. The main arithmetic point in Section~\ref{sec:bounds} is
that taking a power multiplies the initial valuation of a local series,
while leaving its loss per Taylor coefficient unchanged. This permits
the original denominator and decay estimates to scale together.

Small forms could nevertheless be zero. To resolve that obstacle,
Section~\ref{sec:dilation} follows the prime-dilation strategy of
\cite[Section 6]{shifted}, which uses the auxiliary-prime principle
of Lai and Sprang \cite[Lemma 2.3]{laisprang}. At indices equal to a fixed base
index times a large prime, a rescaling makes all beta coefficients
vanish modulo the prime, while the rational constant reduces to a fixed
seed constant. Our adaptations account for Krisman's nonuniform pole
orders and for the loss of the two outer zeros under dilation. The
nonzero seed constants are checked in Section~\ref{sec:seed}.

Thus the kernel and its estimates come from \cite[Sections 2--6]{krisman};
the nonvanishing strategy follows \cite[Section 6]{shifted}. The transfer
under powers and the dilation calculation for this kernel are the
additional arguments given here. The proof uses no irrationality result
for Catalan's constant.

\subsection*{Prerequisites and notation}
The intended reader is familiar with partial fractions, formal Taylor
series, arithmetic modulo primes, elementary prime valuations, and
asymptotic estimates. Write
$v_p$ for the valuation on $\Q$ normalized by $v_p(p)=1$, with
$v_p(0)=+\infty$. The ring $\Z_{(p)}$ consists of rationals whose
denominators are prime to $p$; an element of this ring is a unit when its
numerator is also prime to $p$. Congruences between rationals are taken
in this ring. These definitions suffice for all the local arguments
below. The prime-number-theorem estimate we use is cited explicitly.

\section{Raising the first beta index}\label{sec:construction}
\subsection{The kernel we retain}
We use Krisman's kernel because its denominator cost is already smaller
than its real decay rate. Its integer zeros will remove initial summands
and control rational correction terms; its half-integer poles produce
beta values. The reflection symmetry will remove odd beta indices.

Fix $r=7/6$ and $N=480$, and let $n$ tend to infinity through positive
multiples of $N$. The choice $480=\operatorname{lcm}(160,6)$ makes all
block endpoints and $rn$ integral, and makes $n$ even. The $28$ blocks,
with total multiplicity $S=35$, have parameters
\[
 (\delta,m)=(0,8),\qquad (\kappa/160,1)
 \quad\left(\kappa\in\{1,\ldots,8\}\cup\{10,12,\ldots,46\}\right).
\]
The pole with address $j$ is at $t=-j-1/2$.
A block $b$ adds multiplicity $m_b$ to each pole address in
$I_b(n)=[\delta_b n,(1-\delta_b)n]\cap\Z$. Thus the total pole order at
address $j$ is
\[
 B_j(n)=\sum_bm_b\ind_{j\in I_b(n)}.
\]
For example, at $n=480$ the outer block contributes eight copies of every
address from $0$ to $480$; the blocks with $\kappa=1$ and $46$ contribute
one copy on $[3,477]$ and $[138,342]$, respectively. The orders increase
from $8$ near the ends to $35$ in the middle, and $B_{n-j}=B_j$.
We often omit the argument $n$ from $B_j(n)$.

Set $\Lambda_n=[-(1+r)n-1,rn]\cap\Z$. The kernel is
\begin{equation}\label{eq:R}
 R_n(t)=\frac{(2t+n+1)\prod_{\ell\in\Lambda_n}(t-\ell)^3}
 {\prod_{j=0}^n(t+j+1/2)^{B_j}}.
\end{equation}
The linear factor $2t+n+1$ corrects the sign of reflection about
$-(n+1)/2$, giving
\begin{equation}\label{eq:reflection}
 R_n(-n-1-t)=R_n(t).
\end{equation}
At the central half-integer pole this factor cancels one order. We keep
the nominal order $B_{n/2}=35$ in the formulas, allowing a zero leading
coefficient.

The excess of pole degree over zero degree grows at rate
\[
 G=\sum_bm_b(1-2\delta_b)-3(1+2r)=\frac{179}{10}.
\]
Indeed,
\[
 3\#\Lambda_n=10n+6,\qquad
 \sum_{j=0}^nB_j=\frac{279}{10}n+35,
\]
so $\deg R_n=-Gn-28$. In particular the kernel is proper and decays
rapidly at infinity. These are the parameters and identities of
\cite[Section 2]{krisman}.

\subsection{The square case, then the general form}
Rivoal \cite[Section 2, Lemma 1]{rivoal} used differentiation to remove
$\zeta(3)$ from linear forms in odd zeta values. Krisman
\cite[Section 1.2]{krisman} adapted this device to even beta values.
We use fourth derivatives instead of second derivatives to remove
$\beta(4)$. For example, a partial-fraction term transforms as
\[
 \frac{c}{(t+j+1/2)^2}
 \quad\xrightarrow{\ d^4/dt^4\ }\quad
 \frac{120c}{(t+j+1/2)^6}.
\]
Alternating summation of the latter term produces $\beta(6)$ and a
rational correction. A simple-pole term would instead produce
$\beta(5)$, but reflection cancels that odd-index contribution.

Fourth differentiation would destroy the triple zeros of $R_n$. Squaring
first gives zeros of order six, so the differentiated function still
vanishes at all the relevant integers. The cost is that the largest
nominal pole order doubles from $35$ to $70$, allowing indices as large
as $70+4=74$. Cubing gives zeros of order nine; six derivatives preserve
them, and the largest possible index is $105+6=111$. Reflection cancels
the $\beta(111)$ term, leaving $110$ as the largest possible index.
This explains the endpoints in the theorem.

For the shared proof, fix an integer $q\ge1$ independently of $n$ and put
\begin{equation}\label{eq:L}
 L_{q,n}=\sum_{t=0}^{\infty}(-1)^t
                  \frac{d^{2q}}{dt^{2q}}\bigl(R_n(t)^q\bigr).
\end{equation}
The zeros now have order $3q>2q$, and the degree bound gives absolute
convergence. Write the partial fractions as
\[
 R_n(t)^q=\sum_{j=0}^n\sum_{i=1}^{qB_j}
                  \frac{a_{q,n;j,i}}{(t+j+1/2)^i}.
\]
The rational correction at address $j$ is the omitted initial part of
the beta series. More precisely, for $s>1$ define
\[
 \sigma_j(s)=\sum_{m=0}^{j-1}\frac{(-1)^m}{(2m+1)^s},
 \qquad
 \sum_{t=0}^{\infty}\frac{(-1)^t}{(t+j+1/2)^s}
       =(-1)^j2^s\bigl(\beta(s)-\sigma_j(s)\bigr).
\]
The identity follows by replacing $t+j$ with a new summation index.
With $(i)_d=i(i+1)\cdots(i+d-1)$, it gives
\begin{equation}\label{eq:expansion}
 L_{q,n}=\sum_{j,i}(-1)^j2^{i+2q}(i)_{2q}a_{q,n;j,i}
                 \bigl(\beta(i+2q)-\sigma_j(i+2q)\bigr).
\end{equation}
Reflection \eqref{eq:reflection} gives
$a_{q,n;n-j,i}=(-1)^ia_{q,n;j,i}$. Since $n$ is even, pairing addresses
$j$ and $n-j$ cancels the beta coefficient whenever $i$ is odd. Thus
\begin{equation}\label{eq:target}
 L_{q,n}\in\Q+\sum_{j=q+1}^{\lfloor37q/2\rfloor}\Q\beta(2j).
\end{equation}
We have obtained exactly the required range of possible indices. It
remains to show that clearing the rational coefficients is inexpensive
enough to preserve decay, and then to prove nonvanishing.

\section{Clearing denominators while preserving decay}\label{sec:bounds}
\subsection{Why the arithmetic estimate survives a power}
The relevant arithmetic is local at each pole. Remove the nominal pole
at $t=-j-1/2$ and use a local variable $z$; the remaining Taylor series is
\[
 f_j(z)=(n-2j+2z)\prod_{\ell\in\Lambda_n}(-j-1/2-\ell+z)^3
                   \prod_{h\ne j}(h-j+z)^{-B_h}.
\]
For the $q$th power, the coefficient of pole order $i$ is therefore
$a_{q,n;j,i}=[z^{qB_j-i}]f_j(z)^q$.

Suppose a series $f(z)=\sum f_hz^h$ satisfies
$v_p(f_h)\ge\nu-\lambda h$, where $\nu\in\Z$ and
$\lambda\in\Z_{\ge0}$. Call this a $(\nu,\lambda)$ bound: $\nu$
is the initial valuation and $\lambda$ bounds the loss for each extra
Taylor degree. The square case displays the useful feature:
\[
 [z^h]f(z)^2=\sum_{u+v=h}f_uf_v,
 \qquad v_p(f_uf_v)\ge2\nu-\lambda(u+v)=2\nu-\lambda h.
\]
The same convolution argument shows that $f^q$ has a $(q\nu,\lambda)$
bound. The initial valuation is multiplied by $q$, but the loss per
degree is unchanged. Consequently
\begin{equation}\label{eq:powerlocal}
 v_p(a_{q,n;j,i})\ge q\nu-\lambda(qB_j-i).
\end{equation}
If a harmonic correction of exponent $i+2q$ costs at most
$(i+2q)\lambda$ powers of $p$, the combined bound is
\begin{equation}\label{eq:costcancel}
 q\nu-\lambda(qB_j-i)-\lambda(i+2q)
       =q\bigl(\nu-\lambda(B_j+2)\bigr).
\end{equation}
This is exactly $q$ times the bound for the original second derivative.
It explains both the choice of derivative order $2q$ and the use of the
$q$th power of the original denominator multiplier.

\subsection{The multiplier and its prime bounds}
We now specify that multiplier so that every prime range can be checked.
The notation below records the counts of numerator and denominator
factors divisible by a prime, or their full valuations when prime powers
must be counted. Put
\[
 M_n=19n/3+3,\quad Z=3\#\Lambda_n=10n+6,\quad
 \lambda_p=\lfloor\log_pM_n\rfloor\ (p\text{ odd}),\quad
 \lambda_2=\lfloor\log_2n\rfloor.
\]
All the nonzero integer differences in these local factors have
absolute value below $M_n$. For an odd prime $p>n$, none of the pole
differences is divisible by $p$. For odd $p\le n$, the condition
$p^2>M_n$ means that each nonzero difference contains at most one factor
of $p$, so divisibility counts suffice. Smaller primes can occur to
higher powers and require full valuations. The prime $2$ needs separate
treatment because the pole locations are half-integers.

For $w_{j,\ell}=2j+1+2\ell$, define
\begin{align*}
 \epsilon_p(j)&=3\#\{\ell\in\Lambda_n:p\mid w_{j,\ell}\}
 +\ind_{p\mid n-2j}-\sum_{h\ne j}B_h\ind_{p\mid h-j},\\
 \nu_p(j)&=3\sum_{\ell\in\Lambda_n}v_p(w_{j,\ell})
             -\sum_{h\ne j}B_hv_p(h-j)\quad(p\text{ odd}),\\
 \nu_2(j)&=1-Z-\sum_{h\ne j}B_hv_2(h-j).
\end{align*}
The factor $3$ counts the triple zeros. At $2$, each half-integer
numerator factor costs one power of $2$, explaining $-Z$, while the
linear factor contributes the initial $+1$.

Krisman's full-series estimate \cite[Lemma 4.3]{krisman}, with the choices
in the proof of his Corollary 4.4, gives the following bounds for $f_j$:
\begin{center}
\begin{tabular}{ll}\toprule
Prime range & $(\nu,\lambda)$ bound\\\midrule
$p>n$, odd & $(0,0)$\\
$p\le n$, odd, $p^2>M_n$ & $(\epsilon_p(j),1)$\\
$p\le n$, odd, $p^2\le M_n$ & $(\nu_p(j),\lambda_p)$\\
$p=2$ & $(\nu_2(j),\lambda_2)$\\\bottomrule
\end{tabular}
\end{center}
These are bounds on the whole local series, obtained by multiplying
linear numerator factors and geometric denominator series. We can
therefore apply \eqref{eq:powerlocal} to them.

The original denominator multiplier is
\begin{equation}\label{eq:E}
 E_p(n)=\begin{cases}
 \max_j(-\nu_2(j)+35\lambda_2),&p=2,\\
 0,&p>n\text{ odd},\\
 \max_j(B_j+2-\epsilon_p(j)),&p\le n\text{ odd},\ p^2>M_n,\\
 \max_j(-\nu_p(j)+37\lambda_p),&p\le n\text{ odd},\ p^2\le M_n,
 \end{cases}
 \qquad D_n=\prod_{p\le n}p^{E_p(n)}.
\end{equation}
Here $35$ is the maximum nominal pole order and $37=35+2$ includes
the two derivatives in the original form. Some exponents $E_p$ may be
negative, so $D_n$ is a positive rational multiplier rather than
necessarily an integer. What matters is the integrality of the
coefficients after multiplication.

\begin{proposition}[Arithmetic transfer]\label{prop:arithmetic}
Every rational coefficient of \eqref{eq:expansion}, after collecting equal
beta values and the constant term, becomes an integer on multiplication
by $D_n^q$.
\end{proposition}
\begin{proof}
First take an odd prime $p\le n$. Each odd denominator in
$\sigma_j(i+2q)$ has valuation at most the $\lambda$ in the relevant
row of the table. Equations \eqref{eq:powerlocal}--\eqref{eq:costcancel}
therefore bound the valuation of a constant summand below by $-qE_p(n)$.
The same lower bound holds without the harmonic factor, hence for every
beta coefficient. The factors $(i)_{2q}$ and $2^{i+2q}$ introduce no
odd-prime denominators.

At $2$ the harmonic denominators are odd. Keeping the factor $2^{i+2q}$,
the valuation of a summand is at least
\[
 q\nu_2(j)-\lambda_2(qB_j-i)+i+2q
       \ge q\nu_2(j)-35q\lambda_2\ge-qE_2(n).
\]

It remains to treat odd primes $p>n$. All local coefficients are then
$p$-integral. A harmonic denominator can still be divisible by $p$, so
we use the integer zeros to replace a possibly bad partial sum by its
complement. The use of zeros throughout the pole region to control the
rational constant already appears in Zudilin
\cite[proof of Lemma 1]{zudilin}. We follow Krisman's derivative version
\cite[Lemma 4.6]{krisman}, with zero order $3q$ and derivative order $2q$.
For $1\le\ell\le n$, put
\[
 T(\ell,j)=\sum_i(i)_{2q}a_{q,n;j,i}(j-\ell+1/2)^{-i-2q}.
\]
Reindexing the harmonic sums expresses the rational constant as
$-\sum_{\ell=1}^n(-1)^\ell\sum_{j=\ell}^nT(\ell,j)$.
Because $-\ell\in\Lambda_n$ and $3q>2q$, evaluation of the full
partial-fraction sum gives
\[
 \sum_{j=0}^nT(\ell,j)=(R_n^q)^{(2q)}(-\ell)=0,
 \qquad
 \sum_{j=\ell}^nT(\ell,j)=-\sum_{j=0}^{\ell-1}T(\ell,j).
\]
The prime $p$ cannot divide denominators $2(j-\ell)+1$ on both sides:
two such addresses would differ by a nonzero multiple of $p$ of absolute
value at most $n<p$. Thus one side is a sum of $p$-integral terms, and
the identity makes both sides $p$-integral. This proves the result at
every prime.
\end{proof}

\subsection{The remaining exponential margin}
We next compare the size of $D_n^q$ with the decay of the differentiated
kernel. Estimating logarithms of the zero and pole products by integrals
leads to the profile $\Phi$ below: it records the part of $\log R_n$
proportional to $n$ after the factorial term has been removed. The
constants $A$ and $V$ are the certified odd-prime cost and an upper bound
on this profile, respectively. Define $\mathcal F(y)=y\log y$ for $y>0$
and $\mathcal F(0)=0$, and set
\begin{align*}
 \Phi(x)&=3\{\mathcal F(x+1+r)-\mathcal F(x-r)\}\\
 &\hspace{1cm}-\sum_bm_b\{\mathcal F(x+1-\delta_b)-\mathcal F(x+\delta_b)\},
                     \qquad x\ge r,\\
 A&=\frac{12507029707866288048422972493354331}
 {790300921079786767495600904601600},\qquad V=-\frac{744}{25}.
\end{align*}
The bounds we use from \cite[Corollary 5.9 and Propositions 6.4--6.5]{krisman}
are
\begin{align}
 \log D_n&\le G\log(n!)+n(A+20\log2+o(1)),\label{eq:Dsize}\\
 0<R_n(t)&\le\exp\{-G\log(n!)+n\Phi(t/n)+O(\log n)\}
                              &&(t\ge rn+1),\label{eq:Rsize}\\
 \Phi(x)&\le V&&(x\ge r),\label{eq:phase}\\
 \Phi(x)&\le V-4\log(x/3)&&(x\ge3).\label{eq:tail}
\end{align}
The error in \eqref{eq:Rsize} is uniform in $t$.
The factorial terms in \eqref{eq:Dsize} and \eqref{eq:Rsize} cancel.
What remains is the odd-prime cost $A$, the dyadic cost $20\log2$, and
the real decay exponent $V$. Their sum satisfies
\begin{equation}\label{eq:negative}
 Q:=V+A+20\log2<-0.0714<0.
\end{equation}
This strict margin is the reason for retaining this particular kernel.
Taking a fixed power multiplies both the cost and the decay by the same
number, so it preserves the margin.

\begin{proposition}[Analytic transfer]\label{prop:decay}
For fixed $q$, $D_n^qL_{q,n}\longrightarrow0$ as $n\longrightarrow\infty$
through multiples of $480$. More precisely,
\[
 |D_n^qL_{q,n}|\le\exp\{qnQ+o(n)\}.
\]
\end{proposition}
\begin{proof}
All summands with $0\le t\le rn$ vanish because the zeros have order
$3q$. On the remaining real axis, distances to the zeros and half-integer
poles in the logarithmic derivatives of $R_n$ are at least $1$.
For each fixed $h\ge1$,
\[
 |(\log R_n)^{(h)}(t)|\le(h-1)!(1+Z+\sum B_j)=O_h(n).
\]
The formula for derivatives of an exponential therefore gives
\[
 |(R_n^q)^{(2q)}(t)|\le C_qn^{2q}R_n(t)^q.
\]
This polynomial factor does not affect the exponential rate.
Equations \eqref{eq:Rsize}--\eqref{eq:phase} bound the $O(n)$ summands
with $rn<t\le3n$. Above $3n$, \eqref{eq:tail} supplies the summable factor
$(3n/t)^{4qn}$, since
\[
 \sum_{t>3n}(3n/t)^{4qn}
       \le\int_{3n}^{\infty}(3n/x)^{4qn}\,dx
       =\frac{3n}{4qn-1}.
\]
Combining these estimates yields
$|L_{q,n}|\le\exp\{-qG\log(n!)+qnV+O_q(\log n)\}$.
Now \eqref{eq:Dsize} and \eqref{eq:negative} give the assertion.
\end{proof}

Under joint rationality, the cleared forms would now give integers
tending to zero, so they would eventually have to vanish. The rest of
the proof excludes that possibility.

\section{A fixed seed detects nonzero forms}\label{sec:dilation}
\subsection{Why we use prime indices and integer coordinates}
The auxiliary-prime criterion of Lai and Sprang
\cite[Lemma 2.3]{laisprang} gives a nonvanishing test under the rationality
hypothesis. Outside the finitely many primes
dividing their denominators, rational beta values are $p$-integral.
If a rescaled form has all its beta coefficients divisible by $p$, but
has a rational constant that is a unit modulo $p$, then that form cannot
be zero. Here we only test nonvanishing: this local rescaling need not
have the small real size of $D_n^qL_{q,n}$.

We apply the test at $n=Np$, with the base index $N=480$ fixed and the odd
prime $p$ increasing. Every block endpoint is then a multiple of $p$,
so the poles whose addresses are multiples of $p$ reproduce the base
pole pattern. Following \cite[Section 6]{shifted}, we separate those
factors from the factors that are units modulo $p$.

To do this with integer coordinates, put $x=2t+n+1$. The half-integer
poles become the even integers $n-2a$, and the integer zeros become odd
integers. Define
\begin{equation}\label{eq:F}
 F_n(x)=\frac{x\prod_{\substack{z\text{ odd}\\|z|\le(1+2r)n+1}}(x-z)^3}
 {\prod_{a=0}^n(x-n+2a)^{B_a(n)}}.
\end{equation}
Counting the factors of $2$ gives the exact relation
\begin{equation}\label{eq:change}
 R_n(t)=2^{Gn+29}F_n(2t+n+1).
\end{equation}
Let $c_{q,n;a,i}$ be the coefficient of $(x-n+2a)^{-i}$ in the partial
fractions of $F_n^q$. Define the corresponding form in these coordinates by
\[
 \ell_{q,n}=\sum_{a=0}^n\sum_{i=1}^{qB_a(n)}(-1)^a(i)_{2q}c_{q,n;a,i}
                      \bigl(\beta(i+2q)-\sigma_a(i+2q)\bigr).
\]
It is the alternating sum of the $2q$th derivative with respect to $x$,
evaluated at $x=2t+n+1$. Thus
\begin{equation}\label{eq:ellL}
 L_{q,n}=2^{q(Gn+29)+2q}\ell_{q,n}.
\end{equation}
The power of $2$ is a unit at every odd prime, and the two forms vanish
at exactly the same indices. Denote the rational constant of $\ell_{q,n}$ by
\begin{equation}\label{eq:rho}
 \rho_{q,n}=-\sum_{a=0}^n\sum_{i=1}^{qB_a(n)}
                   (-1)^a(i)_{2q}c_{q,n;a,i}\sigma_a(i+2q).
\end{equation}

\subsection{Which seed survives dilation}
At the base index put
\[
 L=(1+2r)N=1600,\qquad
 W=\sum_bm_b(1-2\delta_b)N=13392.
\]
Here $L$ measures the length of the zero pattern before its two endpoints
are included, while $W$ sums the block widths with multiplicity.
The full pole count at $N$ is $W+S$, with $S=35$.
We may take $p>N+L+1$ throughout this section; in particular $p>2N$.

After substituting $x=pu$, the odd numerator zeros divisible by $p$ have
the form $pv$ with $v$ odd and $|pv|\le Lp+1$. Since $L$ is even, this
means $|v|\le L-1$. The endpoint loss is visible in a small example:
with $L=4$ and $p=5$, the odd zeros from $-21$ to $21$ that are divisible
by $5$ are $-15,-5,5,15$. Dividing by $5$ leaves $-3,-1,1,3$; it does
not leave the extra endpoints $\pm5$ of the untrimmed base pattern.

For our parameters, the surviving seed is consequently
\begin{equation}\label{eq:trim}
 \widetilde F_N(u)=
 \frac{u\prod_{\substack{z\text{ odd}\\|z|\le1599}}(u-z)^3}
 {\prod_{a=0}^N(u-N+2a)^{B_a(N)}}.
\end{equation}
Thus the endpoints are $\pm1599$ rather than $\pm1601$.
The selected pole addresses are $pa$, and
$B_{pa}(Np)=B_a(N)$. The powers of $p$ in the selected factors are
\[
 \begin{array}{c|c}
 \text{factor}&\text{contribution to the exponent of }p\\\hline
 x=pu & 1\\
 L\text{ selected zeros, each triple} & 3L\\
 W+S\text{ selected pole factors} & -(W+S)
 \end{array}
\]
Their sum is $-A_0$, where
\begin{equation}\label{eq:A0}
 A_0=W-3L+S-1=8626.
\end{equation}
This factor count determines the normalization in the next lemma.

Let $\widetilde c_{q;a,i}$ denote the principal-part coefficients of
$\widetilde F_N^q$ at $u=N-2a$. Its associated rational constant is
\begin{equation}\label{eq:seedrho}
 \widetilde\rho_q=-\sum_{a=0}^N\sum_{i=1}^{qB_a(N)}
                  (-1)^a(i)_{2q}\widetilde c_{q;a,i}\sigma_a(i+2q).
\end{equation}
For fixed $q$, the powered seed $\widetilde F_N^q$ and its rational
constant are fixed independently of the varying prime $p$.

\begin{lemma}[Selected factors]\label{lem:selected}
For every sufficiently large odd prime $p$,
\[
 p^{A_0}F_{Np}(pu)\equiv-\widetilde F_N(u)\pmod p.
\]
Reduction commutes with Taylor extraction after the selected poles are
removed. In particular, for $0\le a\le N$ and $1\le i\le qB_a(N)$,
\begin{equation}\label{eq:localcong}
 p^{qA_0-i}c_{q,Np;pa,i}\equiv(-1)^q\widetilde c_{q;a,i}\pmod p.
\end{equation}
\end{lemma}
\begin{proof}
The selected factors counted above give $p^{-A_0}\widetilde F_N(u)$.
Every unselected factor is a unit whose reduction is independent of $u$.
We first compute their combined unit, and then justify extraction of
individual coefficients.

The odd integers from $-Lp+1$ to $Lp-1$ contain each residue class $L$
times. By Wilson's theorem their nonzero residue product is $(-1)^L=1$.
The two extra endpoints $\pm(Lp+1)$ contribute $-1$, also after cubing.
For a block with seed endpoints $a_b,b_b$, each nonzero pole-address
residue occurs $b_b-a_b$ times. The pole differences contain a factor
of $2$, whose product over a nonzero residue cycle is $2^{p-1}=1$ modulo
$p$. The total denominator unit is therefore $(-1)^W=1$. This proves
the rational-function congruence with unit $-1$.

For coefficient extraction put $u=N-2a+z$. The selected principal part
of $F_{Np}(pu)^q$ is
\[
 \sum_{i=1}^{qB_a(N)}p^{-i}c_{q,Np;pa,i}z^{-i}.
\]
After multiplying by $p^{qA_0}z^{qB_a(N)}$, its coefficient of
$z^{qB_a-i}$ is $p^{qA_0-i}c_{q,Np;pa,i}$.
The other selected pole differences are $2(h-a)$ with
$0<|h-a|\le N$, hence units. Every unselected denominator factor also
has a unit constant term and an integral geometric-series inverse.
The pole-removed series therefore lies in $\Z_{(p)}[[z]]$, and reduction
commutes with coefficient extraction. Terms from other poles are
analytic in $z$ and, after pole removal, start in degree $qB_a$.
They cannot contribute to the coefficient in question. Taking $q$th
powers in the congruence proves \eqref{eq:localcong}.
\end{proof}

\subsection{Why the other poles disappear}
The selected poles reproduce the seed. We must also show that the many
other poles contribute nothing to the final reduction. Their residue
classes have fewer denominator factors, which gives the required extra
valuation.

\begin{lemma}[Nonselected coefficients]\label{lem:other}
For every sufficiently large odd prime $p$, if $0\le a\le Np$,
$p\nmid a$, and $1\le i\le qB_a(Np)$, then
\begin{equation}\label{eq:nonselected}
 v_p(c_{q,Np;a,i})\ge-qGN+i.
\end{equation}
\end{lemma}
\begin{proof}
Our choice $p>N+L+1$ makes all nonzero differences between poles, or
between a pole and a numerator zero, smaller than $p^2$ in absolute
value. Any such difference divisible by $p$ has valuation exactly $1$.
At a nonselected pole the factor $x$ is a unit, and every residue class
contains at least $L$ numerator zeros.

Each block $[a_bp,b_bp]$ contains exactly $b_b-a_b$ addresses in any
nonzero residue class. Thus the denominator count at the pole's residue
class is $W$, including its own pole, and is $W-B_a(Np)$ after that
pole is removed. The local analytic part of $F_{Np}^q$ has the bound
$(3qL-q(W-B_a),1)$ in the terminology of Section~\ref{sec:bounds}.
Extracting degree $qB_a-i$ gives
\[
 3qL-q(W-B_a)-(qB_a-i)=3qL-qW+i=-qGN+i,
\]
since $W-3L=GN$. This proves the bound.
\end{proof}

\subsection{The rational constant survives}
Differentiation changes a pole exponent $i$ to $s=i+2q$. To combine
the coefficient congruence with a harmonic denominator of exponent $s$,
rewrite $qA_0-i$ as $q(A_0+2)-s$. This determines the final scaling:
\begin{equation}\label{eq:tau}
 \tau_q=q(A_0+2)=8628q.
\end{equation}
The scaled coefficient at a selected pole then contains $p^s$. This
kills its beta coefficient modulo $p$, but it can cancel a denominator
$p^s$ in the rational correction. That is how the constant can survive
after every beta coefficient has disappeared.

\begin{proposition}[The surviving constant]\label{prop:survive}
For every sufficiently large odd prime $p$, all beta coefficients of
$p^{\tau_q}\ell_{q,Np}$ belong to $p\Z_{(p)}$, and
\begin{equation}\label{eq:constantcong}
 p^{\tau_q}\rho_{q,Np}\equiv(-1)^q\chi_4(p)\widetilde\rho_q\pmod p,
 \qquad \chi_4(p)=(-1)^{(p-1)/2}.
\end{equation}
\end{proposition}
\begin{proof}
Put $s=i+2q$. At a selected pole, \eqref{eq:localcong} gives
\[
 p^{\tau_q-s}(i)_{2q}c_{q,Np;pa,i}
       \equiv(-1)^q(i)_{2q}\widetilde c_{q;a,i}\pmod p.
\]
After multiplication by $p^{\tau_q}$ the corresponding beta coefficient
has valuation at least $s>0$. At a nonselected pole,
\eqref{eq:nonselected} gives valuation at least $q(S+1)+i>0$.

Now consider the rational constant. Every odd denominator in
$\sigma_a(s)$, for $a\le Np$, is less than $2Np<p^2$, so it costs at
most $s$ powers of $p$. The nonselected contributions therefore have
scaled valuation at least
\[
 \tau_q-qGN+i-s=q(S-1)>0.
\]
They vanish in the reduction. The spare $q(S-1)$ comes from the extra
endpoint factors in the selected pole blocks.

For a selected address $pa$, recall the finite sum
\[
 \sigma_{pa}(s)=\sum_{m=0}^{pa-1}\frac{(-1)^m}{(2m+1)^s}.
\]
Multiplication by $p^s$ kills all terms whose denominators are prime to
$p$. The remaining
terms have $2m+1=p(2v+1)$, so
$m=pv+(p-1)/2$ with $0\le v<a$. Hence
\[
 p^s\sigma_{pa}(s)\equiv\chi_4(p)\sigma_a(s)\pmod p,
 \qquad (-1)^m=\chi_4(p)(-1)^v.
\]
The surviving sum is the shorter seed harmonic sum. Combining this
identity with $(-1)^{pa}=(-1)^a$ and the coefficient congruence in
\eqref{eq:rho} gives exactly \eqref{eq:seedrho}, multiplied by
$(-1)^q\chi_4(p)$. This proves \eqref{eq:constantcong}.
\end{proof}

One fixed rational number, $\widetilde\rho_q$, now controls nonvanishing
at all sufficiently large prime indices under the rationality hypothesis.
It remains to show that this number is nonzero for $q=2$ and $3$.

\section{Checking the fixed rational constant}\label{sec:seed}
We do not need the very large numerator and denominator of
$\widetilde\rho_q$. It suffices to find a prime at which the rational
number is integral and has nonzero residue. The auxiliary primes used
for this finite computation are $P=10007$ and $10009$. These fixed primes
have a different role from the varying prime $p$ in Section~\ref{sec:dilation}.

To calculate the constant in \eqref{eq:seedrho}, we need only finitely
many Taylor coefficients at each of the $481$ seed poles. At the pole
$u_a=N-2a$, remove its nominal pole and truncate:
\begin{equation}\label{eq:seedjet}
 J_a(z)=(u_a+z)^q\prod_{\substack{v\text{ odd}\\|v|\le1599}}
                (u_a-v+z)^{3q}
       \prod_{b\ne a}(u_a-u_b+z)^{-qB_b(N)}
       \pmod {z^{qB_a(N)}}.
\end{equation}
The product over $b$ runs through the other addresses in $\{0,\ldots,N\}$.
Then $\widetilde c_{q;a,i}=[z^{qB_a-i}]J_a$. There is no need to expand
the full numerator and denominator of the seed as global polynomials.

The computations are legitimate reductions of rational numbers.
Pole differences have absolute value at most $960$, and the harmonic
denominators at most $959$. The second algorithm below also inverts
nonzero numerator differences, of absolute value at most $2079$, and
Taylor indices less than $105$. All these numbers are smaller than
either auxiliary prime. Primality is checked by trial division through
the square root. The central factor $z^q$ is removed before any
inversion, so no zero constant is inverted there.

Two implementations give the residues in Table~\ref{tab:seeds}.
Either auxiliary prime suffices to prove each required constant nonzero;
the other provides a further check. The $q=1$ row is a control case.
\begin{table}[ht]
\centering
\begin{tabular}{crr}\toprule
Power $q$ & $\widetilde\rho_q\bmod10007$ & $\widetilde\rho_q\bmod10009$\\\midrule
1 (control) & 6332 & ---\\
2 & 3182 & 5359\\
3 & 3335 & 6971\\\bottomrule
\end{tabular}
\caption{Exact residues of the rational seed constants.}
\label{tab:seeds}
\end{table}

The first implementation multiplies truncated polynomials and inverts
the denominator series in \eqref{eq:seedjet}, using reflection to reuse
opposite poles. The second computes every pole independently. It obtains
constant terms from factorial products and higher coefficients from
logarithmic derivatives, as follows.

At a noncentral pole set $H(z)=J_a(z)$ before truncation. At the central
pole first remove the factor $z^q$, leaving a series $H$ with nonzero
constant term. Write
\[
 H(z)=\sum_{j\ge0}h_jz^j,\qquad
 \frac{H'(z)}{H(z)}=\sum_{m\ge1}d_mz^{m-1}.
\]
The factors in \eqref{eq:seedjet} give
\[
 d_m=(-1)^{m-1}\left(q u_a^{-m}
       +3q\sum_v(u_a-v)^{-m}
       -q\sum_{b\ne a}B_b(u_a-u_b)^{-m}\right),
\]
with the first term omitted when $u_a=0$. Equating coefficients in
$H'=H(H'/H)$ gives the recurrence
\[
 h_j=\frac1j\sum_{m=1}^j d_mh_{j-m}.
\]
The initial value $h_0$ is the quotient of the nonzero constant-factor
products. At the central pole the factor $z^q$ is restored afterwards.
This reproduces the required coefficients of $J_a$ using only rational
operations modulo $P$. Substitution in \eqref{eq:seedrho} gives the
table. Both implementations also agree on every even beta coefficient.

\section{The irrationality contradiction}\label{sec:contradiction}
\begin{proof}[Proof of Theorem~\ref{thm:main}]
Suppose all $35$ numbers in the first set are rational. Choose one fixed
integer $H>0$ clearing their denominators, and take $q=2$.
Proposition~\ref{prop:arithmetic} gives
\[
 H D_n^2L_{2,n}\in\Z\qquad(n\in480\Z_{>0}),
\]
and Proposition~\ref{prop:decay} says these integers tend to zero.

Table~\ref{tab:seeds} proves $\widetilde\rho_2\ne0$. Outside the
finitely many primes dividing $H$ or the numerator or denominator of
this fixed rational number, the beta values are $p$-integral and the
seed constant is a unit. Proposition~\ref{prop:survive} therefore gives,
for every sufficiently large such prime,
\[
 p^{17256}\ell_{2,480p}\equiv
       \chi_4(p)\widetilde\rho_2\not\equiv0\pmod p.
\]
The scaled beta terms vanish modulo $p$, leaving the displayed constant.
Hence $\ell_{2,480p}\ne0$, and \eqref{eq:ellL} gives $L_{2,480p}\ne0$.
The excluded set of primes is fixed, so there are infinitely many such
indices. We obtain nonzero integers $H D_{480p}^2L_{2,480p}$ tending to
zero, a contradiction.

For the second set, take $q=3$ and assume its $52$ beta values rational.
The certificate is $\widetilde\rho_3\equiv3335\pmod{10007}$; the
multiplier is $D_n^3$, the local scaling exponent is $25884$, and the
decay exponent is $3Q<0$. The same argument proves the second assertion.
\end{proof}

\begin{corollary}[General power criterion]\label{cor:general}
For any fixed integer $q\ge1$ for which $\widetilde\rho_q\ne0$, at least
one of
\[
 \beta(2q+2),\ \beta(2q+4),\ \ldots,
 \beta\bigl(2\lfloor37q/2\rfloor\bigr)
\]
is irrational.
\end{corollary}
\begin{proof}
Use the same contradiction with $D_n^q$, $\tau_q=8628q$, and the
target span \eqref{eq:target}.
\end{proof}
The seed condition is checked here for $q=1,2,3$. The first case
recovers Krisman's interval; the other two give Theorem~\ref{thm:main}.
For further powers, the nonzero-seed hypothesis remains a separate
condition rather than a uniform conclusion.

\section{Exact checks and audit scope}\label{sec:audit}
\subsection{Computational checks and reproduction}

The mathematical arguments and the finite checks have different roles.
Propositions~\ref{prop:arithmetic}--\ref{prop:decay} and
Lemmas~\ref{lem:selected}--\ref{lem:other} supply bounds for all the
required indices and primes. The computations check the finite seed
conditions and the numerical certificates in the cited estimates.

The two seed algorithms reproduce Table~\ref{tab:seeds} and all the
associated even beta coefficients. They were developed in the same
research session. One uses finite-field polynomial
operations from \texttt{python-flint}; the other uses only the Python
standard library. Four further examples with exact rational coefficients
check the dilation formulas for uniform and nonuniform blocks, including
both classes of odd primes modulo $4$. These finite examples check the
formulas, while the general dilation argument is the proof in
Section~\ref{sec:dilation}.

A separate exact implementation reconstructs all $62,132$ rational
event points of \cite[Proposition 5.6]{krisman} and reproduces the
upward-rounded sum over the $62,131$ intervals. Rational logarithm
enclosures verify the elementary real-phase inequalities and
$Q<-0.0714$. The prime-number-theorem transfer from that arithmetic
certificate to \eqref{eq:Dsize} is used as proved in
\cite[Proposition 5.8 and Corollary 5.9]{krisman}.

The accompanying bundle contains the source, certificate records,
reproduction commands, environment versions, and pinned provenance for
the two principal proof sources, \cite{krisman} and \cite{shifted}.
Its \texttt{README.txt} gives commands for checking
the seeds and the arithmetic budget without additional Python packages.

\subsection{Independent audit}

An independent AI audit reported on 11 November 2026 found no substantive
mathematical error or gap in the version identified by its hash in the
accompanying audit record. The audit checked the general argument and
independently reconstructed the seed and arithmetic certificates using
fresh code, following preliminary inspection and replay of the supplied
checkers. It was not formal verification; the cited Lean project was not
rebuilt.

\begin{thebibliography}{9}
\raggedright
\bibitem{krisman}
Shawn Krisman, \emph{One of $\beta(4),\beta(6),\ldots,\beta(36)$ is irrational},
2026, \href{https://hexagonmath.org/pdf/2610.00285v2}{Hexagon 2610.00285v2}.
Associated source: \url{https://github.com/telaviv/beta-irrationality}.

\bibitem{laisprang}
Li Lai and Johannes Sprang, \emph{Many $p$-adic odd zeta values are irrational},
version 2, 17 February 2025,
\href{https://arxiv.org/abs/2306.10393v2}{arXiv:2306.10393v2}.

\bibitem{rivoal}
Tanguy Rivoal, \emph{Irrationalit\'e d'au moins un des neuf nombres
$\zeta(5),\zeta(7),\ldots,\zeta(21)$}, Acta Arith. \textbf{103} (2002),
157--167, \href{https://doi.org/10.4064/aa103-2-5}{doi:10.4064/aa103-2-5}.

\bibitem{sorokin}
V.\,N. Sorokin, \emph{On the Zudilin--Rivoal theorem},
Math. Notes \textbf{81} (2007), 817--826,
\href{https://doi.org/10.1134/S000143460705029X}{doi:10.1134/S000143460705029X}.

\bibitem{shifted}
twinheterodox, \emph{Irrational even beta values in shifted intervals},
version 2, 6 November 2026, prepared with Codex,
\href{https://hexagonmath.org/pdf/2611.00160v2}{Hexagon 2611.00160v2}.
Associated source: \url{https://github.com/twinheterodox/beta4}.

\bibitem{zudilin}
Wadim Zudilin, \emph{Arithmetic of Catalan's constant and its relatives},
Abh. Math. Semin. Univ. Hambg. \textbf{89} (2019), 45--53,
\href{https://doi.org/10.1007/s12188-019-00203-w}{doi:10.1007/s12188-019-00203-w}.
\end{thebibliography}
\end{document}
